% SPDX-License-Identifier: CC-BY-NC-ND-4.0
% !TEX root = main.tex

\section{アーベル圏}

\begin{definition} \label{def: IsAbelian}
圏$\mathcal{A}$は次の条件を満たすとき\term{あーべるけん}{アーベル圏}{Abelian category}であるという.
\begin{itemize}
    \item 零対象$Z \in \Ob(\mathcal{A})$が存在する.
    \item 任意の$A,B \in \Ob(\mathcal{A})$に対して直積$(A \times B, \pi)$が存在する.
    \item 任意の$f \in \Mor(\mathcal{A})$に対して$f$の核と余核が存在する.
    \item 任意のモノ射$f \in \Mor(\mathscr{C})$に対してある$g \in \Mor(\mathscr{C})$が存在して$f = \ker(g)$となる.
    \item 任意のエピ射$f \in \Mor(\mathscr{C})$に対してある$g \in \Mor(\mathscr{C})$が存在して$f = \coker(g)$となる.
\end{itemize}
\end{definition}

\begin{proposition}
$\mathcal{A}$をアーベル圏とし$f \in \Mor_{\mathcal{A}}(A,B)$とする.\
$f$がモノ射かつエピ射であるとき$f$は同型射となる.\
\end{proposition}

\begin{proof}
$f$はモノ射なのである$C \in \Ob(\mathcal{A})$と$h \in \Mor_{\mathcal{A}}(B, C)$が存在して$f = \ker(h)$となる.\
よって$h \circ f = 0_{A,C} = 0_{B,C} \circ f$であり$f$はエピ射なので$h = 0_{B,C}$である.\
$h \circ \id_B = 0_{B,C}$なので核の普遍性により$f \circ g = \id_B$を満たす$g \colon B \to A$がただ1つ存在する.\
$f \circ g \circ f = \id_B \circ f = f = f \circ \id_A$であり$f$はモノ射なので$g \circ f = \id_A$である.\
よって$f$は同型射である.\
\end{proof}

\begin{remark}
以下断らない限りアーベル圏は局所的に小さいことを仮定する.
\end{remark}

\begin{proposition}
$\mathcal{A}$をアーベル圏とし$f \in \Mor_{\mathcal{A}}(A, B)$をモノ射とする.\
このとき$(A, f) \simeq (\Ker(\coker(f)), \ker(\coker(f)) = (\Im(f), \im(f))$である.\
\end{proposition}

\begin{proposition}
$\mathcal{A}$をアーベル圏とし$A, B \in \Ob(\mathcal{A})$とする.\
$(A \times B, \pi)$を直積とする.
$\iota_0 \coloneq \begin{pmatrix} \id_A \\ 0_{A,B} \end{pmatrix},\ \iota_1 \coloneq \begin{pmatrix} 0_{B,A} \\ \id_B \end{pmatrix}$とすると$(A \times B, \pi, \iota)$は双積になる.
\end{proposition}

\begin{proof}
$\pi_0 \circ \iota_0 = \id_A$なので$(\pi_0, \iota_0)$はレトラクションセクションペアである.\
$X \in \Ob(\mathcal{A})$として$f \in \Mor_{\mathcal{A}}(X, A \times B)$とする.\
$\pi_1 \circ f = 0_{A,B}$であると仮定する.\
$\iota_0 \circ u = f$となる$u \in \Mor_{\mathcal{A}}(X, A)$が存在したとすると$\pi_0 \circ f = u$である.\
そこで$u \coloneq \pi_0 \circ f$とすると$\pi_0 \circ (\iota_0 \circ u) = u = \pi_0 \circ f$であり$\pi_1 \circ (\iota_0 \circ u) = 0_{A,B} \circ u = 0_{A,B} = \pi_1 \circ f$である.\
したがって直積の普遍性により$\iota_0 \circ u = f$である.\
よって次の図式が可換になり$\Ker(\pi_1) = A,\ \ker(\pi_1) = \iota_0$である.\
% https://q.uiver.app/#q=WzAsNCxbMSwwLCJYIl0sWzAsMiwiQSBcXHRpbWVzIEIiXSxbMiwyLCJCIl0sWzEsMSwiQSJdLFswLDEsImYiLDIseyJjdXJ2ZSI6Mn1dLFsxLDIsIlxccGlfMSIsMl0sWzAsMiwiMF97WCxCfSIsMCx7ImN1cnZlIjotMn1dLFswLDMsIlxccGlfMCBcXGNpcmMgZiIsMV0sWzMsMSwiXFxpb3RhXzAiLDJdLFszLDIsIjBfe0EsQn0iLDFdXQ==
\[\begin{tikzcd}
	& X & \\
	& A \\
	{A \times B} && B
	\arrow["{\pi_0 \circ f}"{description}, from=1-2, to=2-2]
	\arrow["f"', curve={height=12pt}, from=1-2, to=3-1]
	\arrow["{0_{X,B}}", curve={height=-12pt}, from=1-2, to=3-3]
	\arrow["{\iota_0}"', from=2-2, to=3-1]
	\arrow["{0_{A,B}}"{description}, from=2-2, to=3-3]
	\arrow["{\pi_1}"', from=3-1, to=3-3]
\end{tikzcd}\]
同様に$\Ker(\pi_0) = B,\ \ker(\pi_0) = \iota_1$である.\

$\pi_1$は$レトラクション$であるからエピ射である.\
よってある$Y \in \Ob(\mathcal{A})$と$g \in \Mor_{\mathcal{A}}(Y, A \times B)$が存在して$\pi_1 = \coker(g)$となる.\
$X \in \Ob(\mathcal{A})$とし$f \in \Mor_{\mathcal{A}}(A \times B, X)$が$f \circ \iota_0 = 0_{A,X}$を満たすとする.\
$\pi_1 = \coker(g)$なので$\pi_1 \circ g = 0_{Y,B}$である.\
よって核の普遍性により$g = \iota_0 \circ v$となる$v \colon Y \to A$がただ1つ存在する.\
$f \circ \iota_0 = 0_{A,X}$なので$f \circ g = f \circ \iota_0 \circ v = 0_{A,X} \circ v = 0_{Y,X}$となる.\
したがって$\pi_1 = \coker(g)$なので余核の普遍性により$f = \pi_1 \circ u$となる$u \colon B \to X$がただ1つ存在する.\
ゆえに$\pi_1 = \coker(\iota_0)$である.\
同様に$\pi_0 = \coker(\iota_1)$となる.\

\end{proof}

\begin{proposition}
アーベル圏は加法圏である.
\end{proposition}

\begin{proof}
$\mathcal{A}$を局所的に小さいアーベル圏とする.\
任意の$A,B \in \Ob(\mathcal{A})$に対して$\Mor_{\mathcal{A}}(A,B)$がアーベル群となることを示す.\
$(A \oplus A, \pi^A, \iota^A)$と$(B \oplus B, \pi^B, \iota^B)$をそれぞれ双積とする.\
$f,g \in \Mor_{\mathcal{A}}(A,B)$に対して$f + g \coloneq \begin{pmatrix} \id_B & \id_B \end{pmatrix} \circ (f \oplus g) \circ \begin{pmatrix} \id_A \\ \id_A \end{pmatrix}$とする.\

\begin{center}
\begin{tikzpicture}[
    commutative diagrams/every diagram,
    x={(1cm,0cm)}, y={(0.7cm,0.5cm)}, z={(0cm,1cm)}
  ]
  % vertices
  \node (BB) at (0,0,4.6)   {$B \oplus B$};
  \node (BL) at (-5,0,0)    {$B$};
  \node (BR) at (5,0,0)     {$B$};
  \node (A0) at (0,-3.6,0)  {$A$};
  \node (A1) at (0,3.6,0)   {$A$};
  \node (AA) at (0,0,0)     {$A \oplus A$};
 
  \path[commutative diagrams/.cd, every arrow, every label]
    % back corner of the base (drawn first: it passes behind the axis)
    (A1) edge node[swap,pos=.75] {$0_{A,B}$} (BL)
    (A1) edge node[swap,pos=.6] {$0_{A,B}$} (BR)
    (A1) edge node[swap,pos=.5] {$\iota_1^A$} (AA)
    % vertical axis
    (AA) edge[commutative diagrams/crossing over]
      node[commutative diagrams/description, pos=.50, font=\footnotesize]
      {$\begin{pmatrix} \begin{pmatrix} f & 0_{A,B} \end{pmatrix} \\
                        \begin{pmatrix} 0_{A,B} & 0_{A,B} \end{pmatrix} \end{pmatrix}$} (BB)
    % product cone
    (BB) edge node[swap] {$\pi_0^B$} (BL)
    (BB) edge node {$\pi_1^B$} (BR)
    % mediating maps out of the centre (labels on the near side of the base)
    (AA) edge node[pos=.55] {$\begin{pmatrix} f & 0_{A,B} \end{pmatrix}$} (BL)
    (AA) edge node[pos=.3] {$\begin{pmatrix} 0_{A,B} & 0_{A,B} \end{pmatrix}$} (BR)
    % front corner of the base (labels on the outside)
    (A0) edge node {$f$} (BL)
    (A0) edge node[swap] {$0_{A,B}$} (BR)
    (A0) edge node[swap,pos=.5] {$\iota_0^A$} (AA);
\end{tikzpicture}
\end{center}
\end{proof}
